% Source: https://artofproblemsolving.com/community/c3h1723237_2019_mathcounts_handbook_transformations_and_measurement_stretch
% Stable identity: aops-topic-1723237-q02
\textbf{Reflection.} The reflection of $(p,q)$ across $y=x-2$ is $(q+2,\,p-2)$: the midpoint $\left(\dfrac{p+q+2}{2},\dfrac{p+q-2}{2}\right)$ lies on the line, and the segment joining the point to its image has direction $(q+2-p,\,p-2-q)$, which is perpendicular to the line's direction $(1,1)$. So $R'=(3+2,\,-5-2)=(5,-7)$.

\textbf{Rotation.} A $90^\circ$ clockwise rotation about the origin sends $(p,q)$ to $(q,-p)$, so $R''=(-7,-5)$.

\begin{asy}
usepackage("amsmath");
size(7cm);
defaultpen(fontsize(10pt));
pair R=(-5,3),R1=(5,-7),R2=(-7,-5),O=(0,0);
draw((-9,0)--(8,0),gray,Arrow(5));draw((0,-9)--(0,6),gray,Arrow(5));
draw((-4,-6)--(7,5),linewidth(0.8));label("$y=x-2$",(7,5),E);
draw(R--R1,dashed);draw(O--R1,dotted+linewidth(0.8));draw(O--R2,dotted+linewidth(0.8));
draw(R--R2,red+linewidth(0.9));
dot(R);dot(R1);dot(R2);dot(O);
label("$R(-5,3)$",R,W);label("$R'(5,-7)$",R1,E);label("$R''(-7,-5)$",R2,W);label("$O$",O,NE);
label("$2\sqrt{17}$",(R+R2)/2,W,red);
label("\mbox{}",(min(currentpicture,user=true)+max(currentpicture,user=true))/2,black); // keeps the text after the figure black when LaTeX typesets the labels inline
\end{asy}

\textbf{Distance.} $RR''=\sqrt{(-7+5)^2+(-5-3)^2}=\sqrt{4+64}=\sqrt{68}=\boxed{2\sqrt{17}}$.
